% ============================================================
\chapter{Connections and Covariant Derivatives}
\label{ch:connections}
% ============================================================

How do you differentiate a vector field on a sphere? On $\R^n$, you
differentiate component by component. But on a curved manifold, tangent
spaces change from point to point, and comparing two vectors at different
points has no intrinsic meaning. An additional ingredient is needed: a
\emph{connection}, which prescribes how to ``transport'' a vector along a
curve. It was Tullio Levi-Civita who, in 1917, formalized this idea in the
Riemannian setting, showing that there exists a unique connection compatible
with the metric and torsion-free.

% ------------------------------------------------------------
\section{Motivation: directional derivatives are not intrinsic}
% ------------------------------------------------------------

On $\R^n$, the derivative of a vector field $Y$ in the direction $X$ is
$(D_X Y)^i = X^j \frac{\partial Y^i}{\partial x^j}$. On a general manifold,
this operation is not intrinsic: it depends on the choice of coordinates. The
notion of \emph{connection} provides the correct framework for differentiating
sections of vector bundles.

\section{Connections on vector bundles}

\begin{definition}[Linear connection]
Let $E \to M$ be a smooth vector bundle. A \emph{connection} (or \emph{covariant
derivative}) on $E$ is an $\R$-bilinear map:
\[
    \nabla : \mathfrak{X}(M) \times \Gamma(E) \to \Gamma(E),
    \quad (X, s) \mapsto \nabla_X s,
\]
satisfying the following axioms:
\begin{enumerate}
    \item \textbf{$C^\infty(M)$-linearity in $X$:} $\nabla_{fX} s = f\, \nabla_X s$
          for all $f \in C^\infty(M)$.
    \item \textbf{Leibniz rule:} $\nabla_X(fs) = (Xf)\, s + f\, \nabla_X s$
          for all $f \in C^\infty(M)$.
\end{enumerate}
\end{definition}

\begin{remark}
Axiom (1) shows that $\nabla_X s$ at a point $p$ depends only on $X_p$
(not on the entire field $X$). Axiom (2) shows that $\nabla$ is \emph{not}
$C^\infty(M)$-linear in $s$: it is a first-order differential operator.
\end{remark}

\begin{proposition}[Locality]
If $s_1 = s_2$ on an open set $U \subset M$, then $\nabla_X s_1 = \nabla_X s_2$
on $U$. The connection is therefore a local operator.
\end{proposition}

\section{Coordinate expression: Christoffel symbols}

Let $(U, x^1, \ldots, x^n)$ be a local chart and $(e_1, \ldots, e_r)$ a local
frame of $E$ over $U$. The \emph{connection coefficients} $\omega^k_{ij}$ are
defined by:
\[
    \nabla_{\partial_i} e_j = \omega^k_{ij}\, e_k.
\]

For the tangent bundle $E = TM$, these coefficients are denoted $\Gamma^k_{ij}$
and called \emph{Christoffel symbols}:
\[
    \nabla_{\partial_i} \partial_j = \Gamma^k_{ij}\, \partial_k.
\]

\begin{proposition}[Coordinate formula]
If $Y = Y^j \partial_j$ and $X = X^i \partial_i$, then:
\[
    \nabla_X Y = X^i \left(\frac{\partial Y^k}{\partial x^i} + \Gamma^k_{ij} Y^j\right) \partial_k.
\]
In components: $(\nabla_X Y)^k = X^i (\partial_i Y^k + \Gamma^k_{ij} Y^j)
= X^i \nabla_i Y^k$.
\end{proposition}

\begin{proof}
By the Leibniz rule:
\begin{align*}
    \nabla_X Y &= \nabla_{X^i \partial_i}(Y^j \partial_j)
    = X^i \nabla_{\partial_i}(Y^j \partial_j) \\
    &= X^i \left[(\partial_i Y^j)\, \partial_j + Y^j \nabla_{\partial_i}\partial_j\right]
    = X^i \left[(\partial_i Y^j)\, \partial_j + Y^j \Gamma^k_{ij}\, \partial_k\right].
\end{align*}
Relabeling the dummy index: $(\nabla_X Y)^k = X^i(\partial_i Y^k + \Gamma^k_{ij} Y^j)$.
\end{proof}

\section{Change of coordinates}

\begin{proposition}[Transformation of Christoffel symbols]
Under a coordinate change $x^i \mapsto \bar{x}^i$, the Christoffel symbols
transform as:
\[
    \bar{\Gamma}^k_{ij} = \frac{\partial \bar{x}^k}{\partial x^l}
    \frac{\partial x^m}{\partial \bar{x}^i} \frac{\partial x^n}{\partial \bar{x}^j}
    \Gamma^l_{mn}
    + \frac{\partial \bar{x}^k}{\partial x^l}
    \frac{\partial^2 x^l}{\partial \bar{x}^i \partial \bar{x}^j}.
\]
The inhomogeneous term shows that $\Gamma^k_{ij}$ do \emph{not} form a tensor.
\end{proposition}

\begin{intuition}[title=\textbf{Christoffel symbols are not tensors}]
The presence of the term $\frac{\partial^2 x^l}{\partial \bar{x}^i \partial \bar{x}^j}$
reflects the fact that the connection encodes additional information beyond the
manifold structure alone. This ``correction'' is precisely what enables covariant
differentiation.
\end{intuition}

\section{Parallel transport}

\begin{definition}[Parallel field along a curve]
Let $\gamma : [a,b] \to M$ be a smooth curve and $V(t)$ a vector field along
$\gamma$ (i.e.\ $V(t) \in T_{\gamma(t)}M$). We say $V$ is \emph{parallel}
along $\gamma$ if:
\[
    \frac{D V}{dt} := \nabla_{\dot{\gamma}(t)} V = 0.
\]
\end{definition}

In coordinates, the parallelism condition reads:
\[
    \frac{dV^k}{dt} + \Gamma^k_{ij}(\gamma(t))\, \dot{\gamma}^i(t)\, V^j(t) = 0,
    \quad k = 1, \ldots, n.
\]
This is a system of first-order linear ODEs.

\begin{theorem}[Existence and uniqueness of parallel transport]
Let $\gamma : [a,b] \to M$ be a smooth curve and $v_0 \in T_{\gamma(a)}M$.
There exists a unique parallel field $V(t)$ along $\gamma$ with $V(a) = v_0$.
The map:
\[
    P_\gamma^{a \to b} : T_{\gamma(a)}M \to T_{\gamma(b)}M,
    \quad v_0 \mapsto V(b),
\]
is a linear isomorphism called \emph{parallel transport} along $\gamma$.
\end{theorem}

\begin{proof}
Existence and uniqueness follow from the Cauchy--Lipschitz theorem applied to
the linear system above. Linearity of $P_\gamma^{a \to b}$ follows from
linearity of the system. Invertibility is obtained by considering parallel
transport along $\gamma$ traversed in reverse.
\end{proof}

\begin{example}[Parallel transport on $S^2$]
Consider the spherical triangle formed by great circle arcs connecting the
North Pole $N = (0,0,1)$ to $A = (1,0,0)$, then $A$ to $B = (0,1,0)$,
then $B$ back to $N$. Parallel transport of a vector around this triangle
rotates it by $\pi/2$. This angle equals the area of the spherical triangle,
illustrating the Gauss--Bonnet theorem.
\end{example}

\section{Torsion of a connection}

\begin{definition}[Torsion tensor]
The \emph{torsion} of a connection $\nabla$ on $TM$ is the $(1,2)$-tensor:
\[
    T(X, Y) = \nabla_X Y - \nabla_Y X - [X, Y].
\]
\end{definition}

\begin{proposition}
$T$ is indeed a tensor: it is $C^\infty(M)$-bilinear and antisymmetric.
In coordinates: $T^k_{ij} = \Gamma^k_{ij} - \Gamma^k_{ji}$.
\end{proposition}

\begin{proof}
We check $C^\infty(M)$-linearity:
\begin{align*}
    T(fX, Y) &= \nabla_{fX} Y - \nabla_Y(fX) - [fX, Y] \\
    &= f\nabla_X Y - f\nabla_Y X - (Yf)X - f[X,Y] + (Yf)X \\
    &= f\, T(X, Y).
\end{align*}
Antisymmetry is immediate. In coordinates:
$T(\partial_i, \partial_j) = \Gamma^k_{ij}\partial_k - \Gamma^k_{ji}\partial_k$
since $[\partial_i, \partial_j] = 0$.
\end{proof}

\begin{definition}[Torsion-free connection]
A connection is \emph{torsion-free} (or \emph{symmetric}) if $T \equiv 0$,
i.e.: $\nabla_X Y - \nabla_Y X = [X, Y]$ for all $X, Y$.
In coordinates: $\Gamma^k_{ij} = \Gamma^k_{ji}$.
\end{definition}

\section{Holonomy}

\begin{definition}[Holonomy group]
Let $(M, \nabla)$ be a manifold with connection and $p \in M$. The
\emph{holonomy group} of $\nabla$ at $p$ is:
\[
    \operatorname{Hol}_p(\nabla) = \{ P_\gamma : T_pM \to T_pM :
    \gamma \text{ smooth loop at } p \} \subset GL(T_pM).
\]
The \emph{restricted holonomy group} $\operatorname{Hol}^0_p(\nabla)$ is the
subgroup corresponding to contractible loops.
\end{definition}

\begin{theorem}[Ambrose--Singer]
The restricted holonomy group $\operatorname{Hol}^0_p(\nabla)$ is a connected
Lie subgroup of $GL(T_pM)$ whose Lie algebra is generated by elements of the
form $P_\gamma^{-1} \circ R(X, Y) \circ P_\gamma$, where $\gamma$ ranges over
paths from $p$ to a point $q$ and $X, Y \in T_qM$.
\end{theorem}

\begin{example}[Holonomy of $S^n$]
For the sphere $S^n$ ($n \geq 2$), $\operatorname{Hol}(S^n) = SO(n)$, since the
sphere is an irreducible simply connected symmetric space.
\end{example}

\section{Covariant derivative of tensors}

\begin{definition}[Extension to tensors]
Let $\nabla$ be a connection on $TM$. We extend $\nabla$ to all tensors by:
\begin{enumerate}
    \item $\nabla_X f = Xf$ for $f \in C^\infty(M)$.
    \item $\nabla_X(S \otimes T) = (\nabla_X S) \otimes T + S \otimes (\nabla_X T)$
          (Leibniz rule).
    \item $\nabla$ commutes with contractions.
    \item On $1$-forms: $(\nabla_X \omega)(Y) = X(\omega(Y)) - \omega(\nabla_X Y)$.
\end{enumerate}
\end{definition}

\begin{proposition}[Covariant derivative of the metric]
For a symmetric $(0,2)$-tensor $g$:
\[
    (\nabla_X g)(Y, Z) = X(g(Y,Z)) - g(\nabla_X Y, Z) - g(Y, \nabla_X Z).
\]
We say $\nabla$ is \emph{compatible with $g$} if $\nabla g = 0$, i.e.\
$(\nabla_X g)(Y,Z) = 0$ for all $X, Y, Z$.
\end{proposition}

\begin{keyformulas}[title=\textbf{Covariant derivative in coordinates}]
For a tensor of type $(r, s)$ with components $T^{i_1 \cdots i_r}_{j_1 \cdots j_s}$:
\[
    \nabla_k T^{i_1 \cdots i_r}_{j_1 \cdots j_s}
    = \partial_k T^{i_1 \cdots i_r}_{j_1 \cdots j_s}
    + \sum_{m=1}^r \Gamma^{i_m}_{k\ell}\, T^{i_1 \cdots \ell \cdots i_r}_{j_1 \cdots j_s}
    - \sum_{m=1}^s \Gamma^\ell_{k j_m}\, T^{i_1 \cdots i_r}_{j_1 \cdots \ell \cdots j_s}.
\]
\end{keyformulas}

\section{Metric compatibility and consequences}

\begin{proposition}[Characterization of compatibility]
The following conditions are equivalent:
\begin{enumerate}
    \item $\nabla g = 0$.
    \item $X \inner{Y}{Z} = \inner{\nabla_X Y}{Z} + \inner{Y}{\nabla_X Z}$
          for all $X, Y, Z \in \mathfrak{X}(M)$.
    \item Parallel transport preserves the inner product: for any curve $\gamma$
          and all $u, v \in T_{\gamma(a)}M$:
          $\inner{P_\gamma u}{P_\gamma v}_{\gamma(b)} = \inner{u}{v}_{\gamma(a)}$.
    \item In coordinates: $\partial_k g_{ij} = \Gamma^l_{ki} g_{lj} + \Gamma^l_{kj} g_{il}$.
\end{enumerate}
\end{proposition}

\begin{proof}
$(1) \Leftrightarrow (2)$: This is the definition of $\nabla g = 0$.

$(2) \Rightarrow (3)$: If $V(t)$ and $W(t)$ are parallel along $\gamma$, then
$\frac{d}{dt}\inner{V}{W} = \inner{\nabla_{\dot\gamma}V}{W}
+ \inner{V}{\nabla_{\dot\gamma}W} = 0$, so $\inner{V}{W}$ is constant.

$(3) \Rightarrow (2)$: Any vector field along $\gamma$ can be written as a
linear combination of parallel fields, and the condition follows.
\end{proof}

\section{Affine space of connections}

\begin{proposition}
The space of connections on $TM$ is an affine space modeled on
$\Gamma(T^*M \otimes T^*M \otimes TM)$: if $\nabla$ and $\tilde{\nabla}$ are
two connections, then $A(X,Y) = \tilde{\nabla}_X Y - \nabla_X Y$ is a
$(1,2)$-tensor.
\end{proposition}

\begin{proof}
We verify $C^\infty(M)$-bilinearity:
$A(fX, Y) = \tilde{\nabla}_{fX}Y - \nabla_{fX}Y = f(\tilde{\nabla}_X Y - \nabla_X Y) = fA(X,Y)$,
and $A(X, fY) = fA(X,Y) + (Xf)Y - (Xf)Y = fA(X,Y)$.
\end{proof}

\section{Exercises}

\begin{exercise}
Show that parallel transport along a meridian of $S^2$ (from $\theta_0$ to
$\theta_1$, with $\phi$ fixed) sends $\partial_\phi$ to $\partial_\phi$
(appropriately normalized).
\end{exercise}

\begin{exercise}
Compute the Christoffel symbols of the metric $g = dr^2 + r^2 d\theta^2$
in polar coordinates on $\R^2$. Verify that the torsion vanishes.
\end{exercise}

\begin{exercise}
Let $\nabla$ be a metric-compatible connection and $V(t)$ a parallel field
along $\gamma$. Show that $\norm{V(t)}$ is constant.
\end{exercise}

\begin{exercise}
Show that if two connections $\nabla$ and $\tilde{\nabla}$ are both torsion-free
and compatible with the same metric $g$, then $\nabla = \tilde{\nabla}$.
(This is the uniqueness of the Levi-Civita connection, proved in the next chapter.)
\end{exercise}

\begin{exercise}
Let $G$ be a Lie group with left-invariant metric. Show that the connection
defined by $\nabla_X Y = \frac{1}{2}[X, Y]$ for left-invariant $X, Y$ is
torsion-free. When is it metric-compatible?
\end{exercise}

\begin{exercise}
Compute the holonomy group of the flat torus $\mathbb{T}^2$. Compare with
that of $S^2$.
\end{exercise}
